\section{Kolmogorov 方程与矩阵指数}
\label{sec:06-Random-Processes/05-Continuous-Time-Markov-Chains-and-Queues/02}

上一节我们拿到了一件精巧的工具——生成矩阵 \(Q\)。它的每个元素都有清晰的概率解释：非对角元 \(q_{ij}\) 描述从状态 \(i\) 到 \(j\) 的瞬时转移速率，对角元 \(q_{ii}=-q_i\) 是离开 \(i\) 的总速率。有了 \(Q\)，你可以在任意短的时间窗 \(\Delta t\) 内做出预报：\(P_{ij}(\Delta t)\approx q_{ij}\Delta t\)（对 \(i\neq j\)），\(P_{ii}(\Delta t)\approx 1-q_i\Delta t\)。

但现实要的不是下一秒，是半年后、三年后。你的老板不会满足于"如果 \(\Delta t\) 足够小"。

设想上一节结尾那台故障修复机。正常运转时，每年平均故障 2 次（\(\lambda=2\)）；一旦故障，平均每年修复 5 次（\(\mu=5\)）。机器此刻正常。半年后它处于故障状态的概率是多少？

手上只有 \(\Delta t\) 内的近似规则，却要跨越 0.5 年的鸿沟。这一篇就解决这个跨越问题。

\subsection{从小时间拼出大时间——朴素尝试}

最直接的办法是把半年切成很多小块再拼起来。取 \(h=0.5/N\)，让 \(N\) 很大。在每个长度为 \(h\) 的小段里，转移矩阵近似为

\[
P(h)\approx I+hQ.
\]

然后把 \(N\) 个小段串起来：

\[
P(0.5)\approx (I+hQ)^N=\bigl(I+\tfrac{0.5}{N}Q\bigr)^N.
\]

这个形式你肯定见过——\((1+x/n)^n\to e^x\)。直觉上，如果把"矩阵"换成"数"，极限就是 \(e^{0.5Q}\)。

但矩阵不是数。\((I+hQ)^N\) 的极限真的存在吗？即使存在，它是不是等于我们逐个元素取指数？如果直接给 \(N=10000\) 做矩阵乘法，计算量也相当可观——两状态还好，状态多了就是噩梦。

我们需要的不只是一个极限猜想，而是一套微分方程，让 \(P(t)\) 在时间的每个瞬间都服从一个确定的变化规律。

\subsection{从半群性质到微分方程}

短时近似的真正力量不在于直接算乘法，而在于它可以导出导数。对任意 \(t\ge0\) 和 \(h>0\)，在 \(t\) 的基础上再往后走一小段 \(h\)：

\[
P(t+h)=P(t)P(h).
\]

这是 Chapman--Kolmogorov 方程的时间齐次版本：先走到 \(t\)，再从 \(t\) 走到 \(t+h\)，与从 0 直接走到 \(t+h\) 等价。

把 \(P(h)\) 换成小时间展开：

\[
P(t+h)=P(t)(I+hQ+o(h))=P(t)+hP(t)Q+o(h).
\]

移项、除以 \(h\)、令 \(h\downarrow0\)：

\[
\frac{P(t+h)-P(t)}{h}=P(t)Q+o(1)\;\Longrightarrow\;P'(t)=P(t)Q.
\]

这是前向 Kolmogorov 方程。逐分量看：

\[
\frac{d}{dt}P_{ij}(t)=\sum_k P_{ik}(t)q_{kj}.
\]

它固定初始状态 \(i\)，把当前时刻可能处在的所有中间状态 \(k\) 汇总起来，看它们各自以多快的速率流向目标 \(j\)。"前向"的意思是我们在时间的最前端——终点 \(t+h\) 处——追加了一小段演化。

如果行分布 \(\mu(t)=\mu(0)P(t)\)，两端对 \(t\) 求导：

\[
\mu'(t)=\mu(t)Q.
\]

这就是主方程。它用最直白的语言说明了每个状态的概率为什么会变化：变化率 = 从其他状态流入的速率 \(-\) 流向其他状态的速率。

\subsection{后向方程：从初始状态的第一跳出发}

从另一个方向追加小时间同样合法：

\[
P(t+h)=P(h)P(t)=(I+hQ+o(h))P(t)=P(t)+hQP(t)+o(h).
\]

同样的极限手续给出后向 Kolmogorov 方程：

\[
P'(t)=QP(t).
\]

逐分量：

\[
\frac{d}{dt}P_{ij}(t)=\sum_k q_{ik}P_{kj}(t).
\]

与前向的视角不同：后向固定了终点 \(j\)，把初始状态 \(i\) 在第一段极短时间里的可能跳转枚举出来。\(i\) 先跳到某个 \(k\)（按照速率 \(q_{ik}\)），再从 \(k\) 用剩下的时间 \(t\) 奔向 \(j\)。

在时间齐次且状态有限的条件下，\(Q\) 与自身的幂级数当然可交换，因此 \(Q\) 与 \(e^{tQ}\) 也可交换。这意味着前向方程和后向方程给出完全相同的解。但在非齐次过程 \(Q(t)\) 中，不同时刻的 \(Q\) 未必可交换，简单的指数公式就不再成立了——那需要时间排序指数（time-ordered exponential），这是另一个故事。

\subsection{矩阵指数：连续时间的矩阵幂}

方程 \(P'(t)=QP(t)\) 配上初值 \(P(0)=I\)，解的形式是

\[
P(t)=e^{tQ}=I+tQ+\frac{t^2}{2!}Q^2+\frac{t^3}{3!}Q^3+\cdots.
\]

这个定义和标量情形完全平行——把 \(e^{at}\) 的泰勒展开里的 \(a\) 换成矩阵 \(Q\) 即可。离散 Markov 链用 \(P^n\) 推进 \(n\) 步，连续时间 Markov 链用 \(e^{tQ}\) 推进时间 \(t\)。两者都是线性演化半群：一个在离散步上做矩阵幂，一个在连续时间上做矩阵指数。

验证它确实是解：逐项求导，第一项（常数项 \(I\)）消失，后面的每一项指数降一次，正好提出一个 \(Q\)：

\[
\frac{d}{dt}e^{tQ}=Q+\frac{2t}{2!}Q^2+\frac{3t^2}{3!}Q^3+\cdots
=Q\Bigl(I+tQ+\frac{t^2}{2!}Q^2+\cdots\Bigr)=Qe^{tQ}.
\]

于是一阶线性矩阵微分方程的形式解就打通了：生成矩阵 \(Q\) 刻画瞬时演化，矩阵指数把瞬时演化累积成任意有限时间的转移矩阵。

一个容易踩的坑：千万别以为对矩阵取指数就是对每个元素取指数。

\[
e^{tQ}\neq\begin{pmatrix}e^{tq_{ij}}\end{pmatrix}_{ij}.
\]

矩阵乘法的不可交换性深埋在指数级数的每一项里——\(Q^2\) 的对角元不是 \(q_{ii}^2\)，而是 \(\sum_k q_{ik}q_{ki}\)。矩阵的指数和元素的指数之间没有简单的逐点对应。

\subsection{回到那台机器：把数字算出来}

现在可以给老板一个答复了。故障修复机的生成矩阵是

\[
Q=\begin{bmatrix}
-\lambda&\lambda\\
\mu&-\mu
\end{bmatrix}
=\begin{bmatrix}
-2&2\\
5&-5
\end{bmatrix}.
\]

设 \(p(t)=P(X(t)=1\mid X(0)=0)\) 为从正常状态出发、到时刻 \(t\) 处于故障状态的概率。主方程只涉及两个状态，可以手解。写出关于 \(p(t)\) 的标量方程：

\[
p'(t)=\lambda(1-p(t))-\mu p(t)=\lambda-(\lambda+\mu)p(t),
\qquad p(0)=0.
\]

这是一阶线性常微分方程。齐次解是 \(Ce^{-(\lambda+\mu)t}\)，特解是常数 \(\lambda/(\lambda+\mu)\)。匹配初值：

\[
p(t)=\frac{\lambda}{\lambda+\mu}+\Bigl(0-\frac{\lambda}{\lambda+\mu}\Bigr)e^{-(\lambda+\mu)t}.
\]

代入数值：\(\lambda=2,\;\mu=5,\;t=0.5\)：

\[
p(0.5)=\frac{2}{7}-\frac{2}{7}e^{-3.5}
\approx 0.2857-0.2857\times0.0302
\approx 0.2771.
\]

半年后，机器约 27.7\% 的概率处于故障状态。第一项 \(2/7\approx28.6\%\) 是长期故障概率——下一篇会从平稳分布的角度重新得到这个数字。指数项 \(e^{-3.5}\approx0.03\) 衰减得很快：\(\lambda+\mu=7\)，特征时间只有约 \(1/7\approx0.143\) 年（不到两个月），半年后初始条件的影响只剩 3\% 了。

\subsection{Uniformization：用 Poisson 时钟驱动离散链}

矩阵指数在理论上很完美，但数值上直接截断幂级数不一定高效。另一种思路绕回了离散链——它叫 uniformization。

如果有限状态 CTMC 的所有离开速率都有上界，选一个统一速率

\[
\nu\ge\max_i q_i.
\]

用它构造一个带自环的离散转移矩阵：

\[
R=I+\frac{Q}{\nu}.
\]

条件 \(\nu\ge\max_i q_i\) 保证了 \(R\) 的对角元 \(1-q_i/\nu\ge0\)，非对角元 \(q_{ij}/\nu\) 本来非负——\(R\) 确是一个合法的随机矩阵。注意它不是上一节的嵌入跳链：跳链禁止原地不动，这里的 \(R\) 故意保留了自环。

由 \(Q=\nu(R-I)\)，且 \(\nu tI\) 与 \(\nu tR\) 可交换：

\[
\begin{aligned}
P(t)=e^{tQ}=e^{\nu t(R-I)}
=e^{-\nu t}e^{\nu tR}
=e^{-\nu t}\sum_{n=0}^{\infty}\frac{(\nu t)^n}{n!}R^n.
\end{aligned}
\]

这个分解有一个清晰的概率解释。想象一个速率为 \(\nu\) 的 Poisson 时钟在滴滴作响——每次滴答，系统按照 \(R\) 做一次离散转移。其中一部分滴答对应真实的状态跳转，另一部分（自环）是虚跳——系统假装跳了一下，实际留在原地。用 Poisson 概率混合所有可能滴答次数下的 \(n\) 步转移，就得到了连续时间 \(t\) 的真实转移矩阵。

Uniformization 同时完成了两件事：把矩阵指数的数值计算转化为截断 Poisson 分布的求和（每项是一个离散矩阵幂），并且把连续时间和离散时间之间的概念桥梁修得明明白白——Poisson 计数过程是计时的钟，离散 Markov 链是跳转的规则。

\subsection{方向约定的最后提醒}

行分布 \(\mu\) 是 \(1\times n\) 的横向量，演化写作

\[
\mu'(t)=\mu(t)Q,\qquad\mu(t)=\mu(0)e^{tQ}.
\]

如果你习惯竖向量（列分布），转置一下就行：

\[
\mu'(t)=Q^{\trans}\mu(t),\qquad\mu(t)=e^{tQ^{\trans}}\mu(0).
\]

前向/后向的名称在不同教科书里可能对调，取决于作者把 \(P(t)\) 看成左乘初始分布还是右乘初始分布。比记名字更可靠的习惯是：每次见到一个新的方程，回到 \(P(t+h)=P(t)P(h)\) 或 \(P(t+h)=P(h)P(t)\) 的分解方式，自己重新推一遍。三行草稿比一个记忆模糊的"前向是 \(PQ\) 还是 \(QP\)"值得多。

到此，从生成矩阵 \(Q\) 到任意有限时间转移矩阵 \(P(t)=e^{tQ}\) 的通路已经打通。Kolmogorov 前向与后向方程刻画了同一套动力学的两个视角；矩阵指数把瞬时速率累积为有限时间的行为；uniformization 再把连续时间拆回 Poisson 时钟驱动的离散链。下一篇带着这套工具走进长时间极限——平稳分布是否存在、是否唯一，以及为什么流量平衡条件能把它直接写出来。
